\begin{abstract}
Building upon recent advances exemplified by BitNet~\cite{bitnet}, the landscape of Large Language Models (LLMs) is entering a transformative phase characterized by 1-bit model architectures. In this study, we propose a novel 1-bit LLM design—\textbf{\text{BitNet b1.58}}—which utilizes ternary weights \{-1, 0, 1\} for every parameter in the network. This model achieves parity with full-precision (FP16 or BF16) Transformer LLMs of equivalent size and data exposure in terms of perplexity and downstream task effectiveness. Notably, it substantially lowers costs across latency, memory footprint, throughput, and power usage. Crucially, the 1.58-bit LLM establishes an updated scaling law and a framework for training next-generation LLMs that combine strong performance with cost efficiency. Additionally, this innovation introduces a fresh computational paradigm and paves the way for hardware solutions custom-built for 1-bit LLM operations.
\end{abstract}

\section{The Era of 1-bit LLMs}

The advancement of Large Language Models (LLMs) in AI has been marked by their rapid expansion in both size and functionality. While their ability to tackle a diverse spectrum of natural language tasks has been impressive, the burgeoning scale of these models presents significant obstacles for real-world deployment and intensifies worries about sustainability due to escalated resource demands.

To mitigate such challenges, post-training quantization has become a prevalent strategy for generating low-bit inference models~\cite{smoothquant, gptq, quip, quip_sharp}. By lowering the precision of both weights and activations, this approach achieves notable savings in memory usage and computational overhead for LLMs. Progress has moved from conventional 16-bit models towards even more compact 4-bit versions~\cite{gptq, awq}. However, despite its widespread adoption in practical LLM systems, post-training quantization tends to be suboptimal and does not fully realize the efficiency–performance tradeoff that is possible.

Recent advances in 1-bit model architectures, exemplified by BitNet~\cite{bitnet}, offer a compelling pathway for decreasing the computational costs of large language models (LLMs) without sacrificing effectiveness. Traditional LLMs utilize 16-bit floating point representations (FP16 or BF16), with matrix multiplications dominating the overall model structure. As a result, most of the computational burden stems from floating-point additions and multiplications. Conversely, BitNet’s matrix multiplications use only integer addition operations, resulting in substantially reduced energy demands for LLM computation. Since power is a primary constraint on computational throughput in modern chips, this reduction directly enables accelerated inference.

Apart from computation, the cost associated with moving model parameters from DRAM to on-chip accelerator memory (such as SRAM) is significant during inference. While expanding SRAM can improve throughput, it is considerably more expensive than DRAM. One key advantage of 1-bit LLMs lies in their substantially reduced memory usage for both storage and data transfer, decreasing the demands on memory bandwidth. This advantage leads to faster weight loading from DRAM, enhancing inference efficiency and speed.

In this paper, we propose a noteworthy 1-bit LLM variant: \textbf{\text{BitNet b1.58}}, characterized by its ternary parameterization where each weight can take a value from the set \{-1, 0, 1\}. The incorporation of the additional 0 value differentiates this model from the original 1-bit BitNet, effectively yielding a parameter representation with 1.58 bits. \text{BitNet b1.58} retains all the key attributes of its predecessor, notably its specialized computation methodology that almost completely eliminates multiplication operations during matrix multiplication, allowing for highly efficient execution. The model’s energy and memory efficiency remain on par with the original BitNet, while throughput and latency also show marked improvements over FP16 LLM benchmarks. Importantly, \text{BitNet b1.58} introduces two further benefits: first, its capacity for modeling is enhanced by allowing explicit feature filtering due to the presence of 0-valued weights, which notably boosts 1-bit LLM performance; second, our empirical results demonstrate that \text{BitNet b1.58} can achieve parity with full-precision (FP16) models in perplexity and downstream task metrics for models sized at 3B parameters and above, provided identical configuration parameters (such as model size and number of training tokens) are used.

\section{BitNet b1.58}

\text{BitNet b1.58} builds upon the BitNet framework, which is a variant of the Transformer architecture where the standard \emph{nn.Linear} layers are replaced by \emph{BitLinear}. The model is initialized and trained from the ground up, utilizing 1.58-bit quantized weights and 8-bit activations. This iteration includes several updates over the original BitNet, which are outlined as follows.

\paragraph{Quantization Function.} To ensure that weights are constrained to the set $\{-1, 0, +1\}$, an \emph{absmean} quantization scheme is employed. This approach begins by normalizing the weight matrix using its mean absolute value, then rounding each element to the closest value among $\{-1, 0, +1\}$:

\begin{equation}
    \widetilde{W} = \mathrm{RoundClip}(\frac{W}{\gamma + \epsilon}, -1, 1),
\end{equation}

\begin{equation}
    \mathrm{RoundClip}(x, a, b) = \max(a, \min(b, \mathrm{round}(x))),
\end{equation}

\begin{equation}
    \gamma = \frac{1}{nm}\sum_{ij} |W_{ij}|.
\end{equation}

The quantization approach for activations remains consistent with that used in BitNet. However, unlike previous versions, the activations are not normalized to fit within $[0, Q_b]$ before applying non-linearities. Instead, each token’s activations are rescaled to $[-Q_b, Q_b]$ to eliminate the need for zero-point quantization. This adjustment both simplifies the code and facilitates system-level tuning, while any resulting impact on empirical performance has been observed to be minimal.

\paragraph{LLaMA-alike Components.}

LLaMA~\cite{llama, llama2} has become the standard reference architecture among open-source LLMs. To support compatibility and foster collaboration with the open-source community, \text{BitNet b1.58} incorporates several elements characteristic of LLaMA. Specifically, it implements RMSNorm~\cite{rmsnorm}, SwiGLU~\cite{swiglu}, rotary embeddings~\cite{rope}, and eliminates bias terms throughout. This design enables seamless integration of \text{BitNet b1.58} with widely used open-source frameworks such as Huggingface, vLLM~\cite{vllm}, and llama.cpp\footnote{\url{https://github.com/ggerganov/llama.cpp}}, thereby reducing adaptation overhead.

\section{Results}

We conducted a comparative analysis between \text{BitNet b1.58} and our FP16 \text{LLaMA LLM} reproduction across different model sizes. Both sets of models underwent pre-training on the RedPajama dataset~\cite{redpajama} for a total of 100 billion tokens to ensure equivalence in pre-training regimes. To assess zero-shot generalization, we evaluated the models on several linguistic benchmarks: ARC-Easy~\cite{arc}, ARC-Challenge~\cite{arc}, Hellaswag~\cite{hellaswag}, Winogrande~\cite{winoGrande}, PIQA~\cite{piqa}, OpenbookQA~\cite{openbookqa}, and BoolQ~\cite{boolq}. Additionally, validation perplexity metrics were computed on the WikiText2~\cite{wikitext2} and C4~\cite{c4} corpora.

We assessed both the GPU memory consumption and inference latency of \text{LLaMA LLM} and \text{BitNet b1.58}. Measurements utilized the FasterTransformer\footnote{\url{https://github.com/NVIDIA/FasterTransformer}} framework, which provides highly optimized GPU inference for large language models. For \text{BitNet b1.58}, Ladder’s 2-bit kernel~\cite{ladder} was incorporated. We present the inference time per output token, which constitutes the primary computational expense during generation.

Perplexity outcomes and resource usage are summarized in Table~\ref{tab:ppl} for both \text{BitNet b1.58} and \text{LLaMA LLM}. The results indicate that, beginning at the 3B parameter scale, \text{BitNet b1.58} achieves perplexity on par with the full-precision \text{LLaMA LLM}, while delivering a 2.71$\times$ reduction in latency and a 3.55$\times$ decrease in GPU memory utilization. Notably, the 3.9B version of \text{BitNet b1.58} attains inference 2.4 times faster than, and operates at 3.32 times lower memory than, the 3B \text{LLaMA LLM}, while exceeding its performance.

Table~\ref{tab:zero-shot} details the zero-shot task performance. For these evaluations, we adopted the approach implemented in \emph{lm-evaluation-harness}\footnote{\url{https://github.com/EleutherAI/lm-evaluation-harness}}. The empirical results reveal that the discrepancy in performance between \text{BitNet b1.58} and \text{LLaMA LLM} decreases with increasing model scale. Crucially, at the 3B size, \text{BitNet b1.58} achieves parity with the full-precision reference model. Mirroring the perplexity findings, end-task scores demonstrate that \text{BitNet b1.58} 3.9B not only requires less computational resource but also surpasses \text{LLaMA LLM} 3B. Collectively, these results establish \text{BitNet b1.58} as a clear Pareto improvement over contemporary leading large language models.

\paragraph{Memory and Latency}

The model was further expanded to 7B, 13B, and 70B parameters, and we conducted a cost assessment at these scales. As depicted in Figure~\ref{fig:latency-memory-energy}, both latency and memory usage display predictable scaling behaviors, with acceleration gains becoming more pronounced as the model size increases. Notably, at 70B parameters, \text{BitNet b1.58} achieves a speed-up factor of 4.1 over the \text{LLaMA LLM} baseline. This improvement is primarily attributed to the fact that the computational burden of \emph{nn.Linear} intensifies with increasing model size. Memory usage exhibits a parallel pattern, as the embeddings—maintained in full precision—account for a diminished share of total memory in larger architectures. Both latency and memory figures were derived using a 2-bit kernel, indicating that there is additional potential for further cost reductions through future optimizations.

\paragraph{Energy}

We further estimated the energy expenditure of arithmetic operations in both \text{BitNet b1.58} and \text{LLaMA LLM}, with an emphasis on matrix multiplication, which is the predominant contributor to energy usage in LLMs. Figure~\ref{fig:energy} breaks down the energy consumption profile, demonstrating that \text{BitNet b1.58} relies primarily on INT8 addition, whereas \text{LLaMA LLM} involves both FP16 addition and FP16 multiplication. Drawing on the energy model from~\cite{energycost, pokebnn}, our analysis indicates that \text{BitNet b1.58} reduces matrix multiplication energy consumption by a factor of 71.4 compared to the baseline on 7nm technology. We additionally measured end-to-end energy requirements using sequences of 512 tokens. These results highlight that as models grow larger, \text{BitNet b1.58} delivers increasing gains in energy efficiency relative to the FP16 \text{LLaMA LLM}. The reason for this is that the proportion of \emph{nn.Linear} computation expands with model size, while the contributions from other components become relatively minor.

\paragraph{Throughput}

We evaluated the throughput of \text{BitNet b1.58} and the \text{LLaMA LLM} 70B on a system equipped with two 80GB A100 GPUs, utilizing pipeline parallelism~\cite{gpipe} to allow for the execution of the \text{LLaMA LLM} 70B model. The batch size was incrementally increased until the GPUs' memory capacity was reached, with all tests conducted at a sequence length of 512. According to Table~\ref{tab:throughput}, \text{BitNet b1.58} 70B is capable of accommodating batch sizes up to eleven times larger than those supported by \text{LLaMA LLM}, which translates into an 8.9-fold improvement in throughput.

\textbf{\text{BitNet b1.58} establishes a new paradigm for scaling model performance relative to inference costs.} Based on the comparisons presented in Figure~\ref{fig:latency-memory-energy} and \ref{fig:energy}, we can illustrate the equivalency in efficiency across varying model sizes for 1.58-bit and 16-bit models as follows:

\begin{itemize}
    \item BitNet b1.58 at 13B exhibits better efficiency—spanning latency, memory usage, and energy consumption—when compared to a 3B FP16 LLM.
    \item BitNet b1.58 at 30B is more resource-efficient with respect to latency, memory, and energy consumption than a 7B FP16 LLM.
    \item BitNet b1.58 at 70B surpasses a 13B FP16 LLM in terms of latency, memory demands, and energy use.
\end{itemize}

\paragraph{Training with 2T Tokens}

The volume of training tokens plays a significant role in the performance of large language models (LLMs). To investigate the token scalability of \text{BitNet b1.58}, we pre-trained this model variant using 2T tokens, adhering to the same dataset construction guidelines as those employed for StableLM-3B~\cite{StableLM-3B-4E1T}, the leading open-source 3B model. Evaluation was conducted using a suite of benchmarks, including Winogrande~\cite{winoGrande}, PIQA~\cite{piqa}, SciQ~\cite{sciq}, LAMBADA~\cite{lambada}, and ARC-easy~\cite{arc}, with zero-shot accuracy results reported in Table~\ref{tab:2t}. For datasets where both accuracy and normalized accuracy are applicable, their average is computed. The performance metrics for StableLM 3b at 2T tokens were obtained directly from its technical documentation. The data indicate that \text{BitNet b1.58} achieves the highest performance across all evaluated tasks, demonstrating robust generalization abilities characteristic of 1.58-bit LLMs.

\section{Discussion and Future Work}

\textbf{1-bit Mixture-of-Experts (MoE) LLMs}

The Mixture-of-Experts (MoE) framework has demonstrated substantial efficiency for scaling large language models by lowering computational FLOPs. However, MoE models typically suffer from substantial memory demands and heavy inter-chip communication, posing challenges to practical deployment. The introduction of 1.58-bit LLMs offers solutions to these bottlenecks. By minimizing memory usage, fewer hardware devices are needed to deploy MoE architectures. In addition, data transfer requirements for activations between devices are drastically decreased. If the entire model fits within a single chip, communication overhead is effectively eliminated.

\textbf{Native Support for Long Sequence Processing in LLMs}

As LLMs advance, efficient handling of lengthy input sequences has become increasingly essential. A core limitation in this area is the significant memory required for storing key-value (KV) caches during inference. \text{BitNet b1.58} marks a notable advancement in this regard, compressing activations from 16 bits down to 8 bits, which effectively doubles the feasible context window within the same hardware constraints. Furthermore, there exists potential to compress activations to 4 bits or below for 1.58-bit LLMs without incurring quality loss, presenting a promising direction for future exploration.

\textbf{LLMs for Edge and Mobile Environments}

Deploying LLMs with 1.58-bit representations could dramatically boost language model performance on edge and mobile devices, which frequently face tight constraints on memory and computational throughput. Lower memory and energy requirements inherent to 1.58-bit LLMs make it feasible to utilize LLMs on these platforms, paving the way for new applications and enhanced device functionality that were once out of reach. Furthermore, since edge and mobile devices primarily utilize CPUs, and 1.58-bit LLMs such as \text{BitNet b1.58} are well-suited for CPU execution, these models can be run efficiently, unlocking additional opportunities for advanced natural language capabilities.

\textbf{Custom Hardware for 1-bit LLMs}

Emerging solutions like Groq\footnote{\url{https://groq.com/}} have highlighted the advantages and possibilities of constructing tailored hardware—such as LPUs—for LLM acceleration. Looking ahead, we advocate for the research and development of hardware architectures and systems specifically tailored for the demands of 1-bit LLMs, leveraging the computational paradigm presented by BitNet~\cite{bitnet}.

\end{document}